\uuid{Um8A}
\titre{Étude de la classe $\mathcal{C}^1$ de $f(x,y)=x^2y^2\ln(x^2+y^2)$}
\niveau{L2}
\module{Analyse}
\chapitre{Fonctions de plusieurs variables}
\sousChapitre{Continuité et différentiabilité}
\theme{Dérivées partielles, continuité, coordonnées polaires, classe $\mathcal{C}^1$}
\auteur{Maxime Nguyen}
\datecreate{2026-03-24}
\organisation{amscc}
\difficulte{4}

\contenu{
\texte{
  On considère la fonction $f$ définie sur $\mathbb{R}^2$ par
  \[
    f(x,y) = x^2 y^2 \ln(x^2+y^2) \;\text{ si } (x,y) \neq (0,0), \qquad f(0,0) = 0.
  \]
  Cette fonction est-elle de classe $\mathcal{C}^1$ ?
}
\begin{enumerate}
  \item \question{Calculer les dérivées partielles de $f$ pour $(x,y) \neq (0,0)$.}
    \reponse{
      \[
        \frac{\partial f}{\partial x}(x,y) = 2xy^2\ln(x^2+y^2) + \frac{2x^3y^2}{x^2+y^2},
      \]
      \[
        \frac{\partial f}{\partial y}(x,y) = 2x^2y\ln(x^2+y^2) + \frac{2x^2y^3}{x^2+y^2}.
      \]
    }

  \item \question{Calculer les dérivées partielles de $f$ en $(0,0)$ par la définition (taux d'accroissement).}
    \reponse{
      On a $f(h,0) = 0$ et $f(0,h) = 0$ pour $h \neq 0$. Donc $\dfrac{\partial f}{\partial x}(0,0) = \dfrac{\partial f}{\partial y}(0,0) = 0$.
    }

  \item \question{Étudier la continuité des dérivées partielles en $(0,0)$.}
    \reponse{
      En coordonnées polaires $x = r\cos\theta$, $y = r\sin\theta$ :
      \[
        \left|\frac{\partial f}{\partial x}(x,y)\right| \leq 4r^3|\ln r| + 2r^3.
      \]
      Par croissances comparées, $r^3|\ln r| \to 0$ et $r^3 \to 0$ quand $r \to 0$. Donc $\dfrac{\partial f}{\partial x}$ est continue en $(0,0)$. Le raisonnement est identique pour $\dfrac{\partial f}{\partial y}$.

      Ainsi $f$ est de classe $\mathcal{C}^1$ en $(0,0)$, puis $\mathcal{C}^1$ sur $\mathbb{R}^2$ par opérations sur les fonctions $\mathcal{C}^1$.
    }
\end{enumerate}
}
